\documentclass[11pt,a4paper]{article}

% Fix fi/fl ligature extraction from the PDF (correct ToUnicode CMaps)
\input{glyphtounicode}
\pdfgentounicode=1
\usepackage{cmap}
\usepackage[utf8]{inputenc}
\usepackage[T1]{fontenc}
\usepackage{lmodern}
\usepackage{amsmath,amssymb}
\usepackage{geometry}
\usepackage{enumitem}
\usepackage{listings}
\usepackage{xcolor}

\geometry{top=0.4in,bottom=0.4in,left=0.7in,right=0.7in}

\lstset{
  basicstyle=\ttfamily\footnotesize,
  columns=fullflexible,
  keepspaces=true,
  xleftmargin=1.2em,
  aboveskip=0.05em,
  belowskip=0.05em,
  showstringspaces=false,
  language=Haskell
}

\setlength{\parskip}{0pt}
\linespread{0.82}

\usepackage{titlesec}
\titlespacing*{\section}{0pt}{0.7ex plus .2ex}{0.4ex}
\titleformat{\section}{\normalfont\large\bfseries}{\thesection}{0.6em}{}
\titlespacing*{\paragraph}{0pt}{0.35ex plus .1ex}{0.5em}

\begin{document}

\begin{center}
{\Large\bfseries Where a Unit's State Lives}\\[0.2ex]
{\normalsize The Ledger}
\end{center}
\vspace{0.4ex}

%==============================================================================
\section{The Question}
%==============================================================================

A unit is anything that can be held: a future, a bond, a mandate, a strategy
contract. A wallet holds it. A unit's state --- its terms, its lifecycle, a
holder's cost --- lives somewhere. This fixes where. The placement makes one source
of truth, conservation, and deterministic replay attainable. Each economic fact then
has exactly one home and one writer, so internal reconciliation failure cannot arise;
the sealed single-writer discipline, and the purity and totality of replay, make
conservation and deterministic replay hold by construction.

%==============================================================================
\section{The Answer}
\label{sec:answer}
%==============================================================================

Every economic relationship a wallet has that reifies as a single unit it holds --- a
holding, a mandate, a strategy --- has its state in three homes, held in two maps. The
reification, and the scope it bounds, are shown in \S\ref{sec:why}.

One rule governs every placement: each economic fact has exactly one home and one
writer --- single source of truth, so internal reconciliation failure cannot arise. A
fact fails the rule in only two ways: keyed wrong, one home collapses two facts or one
fact duplicates across homes that drift; authored wrong, two authorities keep one fact
and their records drift apart. Each of the two questions forestalls one failure, and
because there are only these two, there is no third question: whether the fact depends
on the holder, and who authors it. \S\ref{sec:why} works each failure on a concrete
unit.

\textit{Does the fact depend on the holder, or only on the unit?} The scope is one
unit's state, so the unit and a holder of it are the only terms in view. By the
reification, every economic fact --- one that enters conservation, valuation, or profit and loss ---
is a (holder, unit) fact or a unit fact: a fact about a relationship is a fact about the
wallet's position in the unit that relationship is, never about a wider key. What stays
keyed by the wallet alone --- KYC, permissions, an audit cursor --- enters none of the
three reckonings; it is identity, not economic state.

\textit{Is the fact's record owned by an outside authority or by the ledger?} The
criterion is who owns the record's history, not who sources the number. An externally
authored fact is one an outside authority --- the exchange, the contract, the
reference-data provider --- owns and restates, and the ledger preserves version by
version. A ledger-authored fact is one the ledger's own events produce and overwrite.

The two questions give a $2\times2$; three of its four cells are occupied. Each
occupied cell is a home: one kind of state.

\begin{center}\footnotesize
\begin{tabular}{@{}r|p{0.39\textwidth}p{0.39\textwidth}@{}}
 & \textbf{ledger-authored} & \textbf{externally authored} \\ \hline
\textbf{per unit} & \textbf{Status}: lifecycle stage, last settlement price, current
weights, benchmark level & \textbf{Terms}: multiplier, currency, expiry, strike, ISIN,
fee schedule, mandate text, benchmark identity \\[0.5ex]
\textbf{per (holder, unit)} & \textbf{Position}: held quantity, high-water mark &
\textit{--- empty ---} \\
\end{tabular}
\end{center}

Each cell names a category; the listings model a slice of it --- for Status, the
lifecycle stage carrying the settlement price, with current weights and benchmark level
elided as the Position row's further fields are. One criterion sorts the whole category: a
benchmark level the ledger records and overwrites is ledger-authored, while the
benchmark identity the provider restates is the authority's versioned record --- the one
externally authored, the other not, though both come from that provider.

The fourth cell --- an externally authored (holder, unit) fact --- is empty
(\S\ref{sec:why}).

Terms and status share the unit key but cannot be one value (\S\ref{sec:why}); they ride
together as a pair --- a third home, not a third map (\S\ref{sec:construction}).

%==============================================================================
\section{Why Three}
\label{sec:why}
%==============================================================================

Each occupied cell exists for one concrete reason; the empty cell is empty for
one.

\paragraph{Position is keyed by (holder, unit) because two holders of one
contract hold different state.}
A buyer and a seller hold the same future. The buyer's held quantity is $+1000$;
the seller's is $-1000$, and their costs differ too. A unit-keyed value would
store one number and collapse the two. The holder must be in the key.

\paragraph{Status is keyed by the unit alone because one value is read
identically by every holder.}
A future settles at one price on the day; an index stands at one level. Stored per
holder, the number is copied across thousands of rows free to drift apart --- a
reconciliation break by construction. Keyed by the unit, it is one value with one
writer.

\paragraph{Terms are a home distinct from status because the two have different
authorities of record.}
Terms are externally authored, status ledger-authored (\S\ref{sec:answer}). Co-mingling
the authority's record with the ledger's own is the single-source-of-truth violation the
system exists to prevent, so the two are distinct values: a non-empty version list grown
by \texttt{appendVersion}, a single value replaced by \texttt{settle}. The amendment
event that grows the list is out of scope here, so within this file every terms value has
exactly one version, and the distinction is exercised only as the contrast between its two
writers.

\paragraph{The fourth cell is empty because no authority issues a fact about one
holder's position.}
A position is the ledger's own from origin: the held quantity folded from the move
stream, a high-water mark written by the valuation event, a write-once entry NAV
(elided here) folded from the opening event. A custodian holding statement or a
prime-broker position report is a reconciliation input the ledger checks against, not a
record it adopts as an authority's own. So the (holder, unit) key carries only
ledger-authored facts, and the externally authored cell is empty by construction.

A managed account looks like the counterexample: its high-water mark and accrued
fee appear to demand a wallet-keyed home. They do not. The manager issues a
mandate unit to the client --- $-1$ to the manager, $+1$ to the client, summing to
zero like any issued unit --- so the client's state under that mandate is a (client,
mandate-unit) position, its high-water mark and entry NAV facts about that position.
Two mandates are two rows; a wallet-keyed value would merge them. This discharges the
reification for one mandate. A relationship spanning several instruments that does not
likewise reduce to a single (holder, unit) row lies outside the scope of
\S\ref{sec:answer}: its fact would be a (holder, several-units) fact occupying a fourth
home and a third map, which the count excludes.

%==============================================================================
\section{The Construction}
\label{sec:construction}
%==============================================================================

The homes are built from the smallest pieces, each forced by the one before. The
Haskell is \texttt{States.hs}; the listings reproduce its declarations, deriving
clauses elided.

\paragraph{A balance is a quantity in exact minor units.}
It has a zero and adds associatively --- a monoid --- and every balance has a
negation, making it a group. The negation later forces the two legs of a transfer
to cancel.

\begin{lstlisting}
newtype Qty = Qty Integer
instance Semigroup Qty where Qty a <> Qty b = Qty (a + b)
instance Monoid    Qty where mempty = Qty 0
negQty :: Qty -> Qty
negQty (Qty n) = Qty (negate n)
\end{lstlisting}

\paragraph{A balance is held by a wallet, of some unit.}
Two keys name it; the holder is in the key (\S\ref{sec:why}). The map keyed by this
pair is assembled below, once its value has
grown past a bare quantity; the \texttt{Maybe} its lookup returns carries a real
distinction, exhibited on the assembled store.

\begin{lstlisting}
newtype WalletId = WalletId String
newtype UnitId   = UnitId   String
\end{lstlisting}

\paragraph{Some facts are one per unit --- shared status.}
Status is one stage of the unit's lifecycle, overwritten
on each settlement. The price a settlement fixes is data the stage carries: it rides
on the \texttt{Active} constructor, so ``active with no price'' and ``registered yet
priced'' are unspellable --- the correlation holds by the type, not by a writer
trusted to set two fields in lockstep. The lifecycle carries only the two stages its
writers reach: \texttt{Registered} at registration --- a unit recorded in and known to
the ledger before any settlement, which an OTC option, a bespoke mandate, or an internal
strategy reaches exactly as an exchange instrument does, not admission to an exchange ---
and \texttt{Active} once the unit
settles. A price is a number but not a quantity --- prices are never added --- so
\texttt{Price} is a separate newtype with neither identity nor inverse, never summed
into a balance.

\begin{lstlisting}
newtype Price = Price Integer
data Lifecycle = Registered | Active Price
newtype UnitStatus = UnitStatus { usLifecycle :: Lifecycle }
defaultStatus :: UnitStatus
defaultStatus = UnitStatus Registered              -- the initial stage, no price
\end{lstlisting}

\paragraph{Terms are the authority's record, grown only by appending.}
``One or more
versions, append-only'' is a non-empty list, so a registered unit always has a
current version and ``registered but versionless'' is not representable. The
constructor is not exported, so no importer can lay down a fresh one-version value
and discard history; \texttt{register}, the one writer that lays down version one,
refuses a unit already present (below). No door shortens history: the
append-only discipline holds by construction, not by convention.

\begin{lstlisting}
data TermsVersion = TermsVersion String           -- multiplier, ISIN, expiry, ...
newtype ProductTerms = ProductTerms (NonEmpty TermsVersion)  -- constructor not exported
currentTerms  :: ProductTerms -> TermsVersion
currentTerms  (ProductTerms vs) = NE.last vs
appendVersion :: TermsVersion -> ProductTerms -> ProductTerms
appendVersion tv (ProductTerms vs) = ProductTerms (vs <> (tv :| []))
\end{lstlisting}

\paragraph{A position carries more than a balance.}
Its primary field, the held quantity \texttt{psBal}, conserves (\S\ref{sec:right}).
Beside it rides a per-(holder, unit) high-water mark: \texttt{psHwm} witnesses that the
Position home can carry a non-conserved fact. Its writer --- a valuation event out of
scope here --- leaves it \texttt{mempty} throughout this file; no move writes it as a
cancelling leg, so it bears no zero-sum invariant, and it keeps the full \texttt{Qty}
type rather than a stripped newtype like \texttt{Price} because its operation is not
settled here --- no aggregate over holders is claimed. A first touch starts a position at
\texttt{zeroP}, any further field elided.

\begin{lstlisting}
data PositionState = PositionState
  { psBal :: Qty       -- held quantity: primary, conserved (sum over holders, sec.5)
  , psHwm :: Qty }     -- high-water mark: not conserved; out-of-scope writer, mempty here
zeroP :: PositionState
zeroP = PositionState mempty mempty
\end{lstlisting}

\paragraph{The three homes, two maps.}
Terms and status share the unit key, so a unit's
entry is the pair (terms, status). Co-presence is then the shape of the map, not an
invariant to police --- no ``in terms but not status'' value can be written, since
one entry carries both halves or neither. The position map is the balance map
enriched, its value carrying the held quantity in \texttt{psBal}. The constructor and
the field selectors \texttt{ledgerUnit} and \texttt{ledgerPS} are not exported, so a
\texttt{Ledger} is built only by the writers below and read only through the accessors:
were a selector exported, record update through it (\texttt{l \{ ledgerPS = ... \}})
would install a non-conserving map without naming the constructor, bypassing the
single-writer discipline. Coherence is carried by the pair, not the seal; the seal is
left to keep conservation true by construction. \texttt{emptyLedger} has both maps empty.

\begin{lstlisting}
data Ledger = Ledger                          -- constructor + field selectors not exported
  { ledgerUnit :: Map UnitId (ProductTerms, UnitStatus) -- terms + status, co-present
  , ledgerPS   :: Map (WalletId, UnitId) PositionState } -- per (holder, unit); psBal
emptyLedger :: Ledger
emptyLedger = Ledger Map.empty Map.empty
\end{lstlisting}

\paragraph{Registration and settlement.}
Registration introduces a unit, writing the pair of its first terms version and
default status as one entry; it refuses a unit already present, returning
\texttt{Nothing}.
Settlement adjusts only the status half --- \texttt{Map.adjust} over the pair's
\texttt{snd} --- overwriting it with the \texttt{Active} stage carrying the day's
price, leaving terms untouched. \texttt{appendVersion}, the remaining terms writer, is
out of scope (\S\ref{sec:why}), so \texttt{register} and \texttt{settle} are the only
writers exercised; neither touches \texttt{psBal}.

\begin{lstlisting}
register :: UnitId -> TermsVersion -> Ledger -> Maybe Ledger
register u tv l
  | Map.member u (ledgerUnit l) = Nothing
  | otherwise = Just (l { ledgerUnit =
      Map.insert u (ProductTerms (tv :| []), defaultStatus) (ledgerUnit l) })
settle :: UnitId -> Price -> Ledger -> Maybe Ledger
settle u px l
  | Map.member u (ledgerUnit l) =
      Just (l { ledgerUnit =
        Map.adjust (\(t, _) -> (t, UnitStatus (Active px))) u (ledgerUnit l) })
  | otherwise = Nothing
\end{lstlisting}

\paragraph{A move on the ledger is a transfer, lifted.}
A move's effect is its per-wallet net delta, written into \texttt{psBal} and gated on
registration. The two legs --- \texttt{negQty q} at the source, \texttt{q} at the
destination, exact inverses in \texttt{Qty} (above) --- are summed per wallet: two
distinct wallets net $-q$ and $+q$; a self-move's legs land on one wallet and cancel to
\texttt{mempty}. Either way the move sums to \texttt{mempty} over wallets.

\begin{lstlisting}
data Move = Move UnitId WalletId WalletId Qty
applyMove :: Move -> Ledger -> Maybe Ledger
applyMove (Move u from to q) l
  | Map.member u (ledgerUnit l) =
      Just (l { ledgerPS = Map.foldrWithKey writeNet (ledgerPS l) (netDeltas from to q) })
  | otherwise = Nothing
  where
    netDeltas f t qty = Map.insertWith (<>) t qty (Map.insertWith (<>) f (negQty qty) Map.empty)
    writeNet w d ps
      | d == mempty = ps                                 -- a zero net writes no row
      | otherwise   =
          let cur = Map.findWithDefault zeroP (w, u) ps
          in  Map.insert (w, u) (cur { psBal = psBal cur <> d }) ps
\end{lstlisting}

The \texttt{Maybe} guards input, not the balance: \texttt{applyMove} returns
\texttt{Nothing} only for a move naming an unregistered unit, never for a balance
that failed to hold --- a well-formed move's two legs cancel by construction, so
\texttt{Nothing} answers ``is this unit known?'', never ``did the balance hold?''. A
zero-quantity move and a self-move ($\texttt{from} = \texttt{to}$) are both well-formed
and accepted, but each nets to \texttt{mempty} on every wallet (above), so neither
writes a row: \emph{held} means named in a move that nets nonzero on it.

\paragraph{Reading a position.}
\texttt{position} looks up the pair key on \texttt{ledgerPS}, the store the system
keeps. Its \texttt{Maybe} carries the distinction. \texttt{Nothing} is \emph{never
held} --- no key. \texttt{Just} a row is \emph{held}, its \texttt{psBal} the held
quantity; a closed-out holding is \texttt{Just} a row with \texttt{psBal} zero ---
\emph{held, and now flat} --- because \texttt{applyMove}'s first touch writes a row
and a close-out leaves it at zero, never deleted. Settlement entitlement and
wash-sale lookback answer \emph{never held} and \emph{held and flat} differently; the
retained row keeps them apart.

\begin{lstlisting}
position :: Ledger -> WalletId -> UnitId -> Maybe PositionState
position l w u = Map.lookup (w, u) (ledgerPS l)
\end{lstlisting}

%==============================================================================
\section{Why It Is Right}
\label{sec:right}
%==============================================================================

\paragraph{Conservation.}
Conservation is an invariant of the writer, not the store type, which can hold a
non-conserving assignment. \texttt{applyMove} is the only function that touches
\texttt{psBal}, and it writes the two legs together from one quantity, so each move
changes the holding sum by $\texttt{negQty q} \mathbin{\texttt{<>}} \texttt{q} =
\texttt{mempty}$. \texttt{register} and \texttt{settle} touch only
\texttt{ledgerUnit}, never \texttt{psBal}, leaving every holding sum unchanged. From
\texttt{emptyLedger}, whose sum is zero, every event preserves it, so every
reachable ledger conserves --- the sealed constructor and withheld field selectors leave
no other door.

\begin{lstlisting}
netBal :: Ledger -> UnitId -> Qty
netBal l u = foldMap psBal [ p | ((_, u'), p) <- Map.toList (ledgerPS l), u' == u ]
\end{lstlisting}

\paragraph{Deterministic replay.}
\texttt{apply} dispatches over the three events; \texttt{replay} folds.

\begin{lstlisting}
data Event = Registered UnitId TermsVersion | Moved Move | Settled UnitId Price
apply :: Event -> Ledger -> Maybe Ledger
apply (Registered u tv) = register u tv
apply (Moved m)         = applyMove m
apply (Settled u px)    = settle u px
replay :: [Event] -> Ledger -> Maybe Ledger
replay events l0 = foldM (flip apply) l0 events
\end{lstlisting}

Replay is deterministic because \texttt{apply} is a pure, total function of the
event and the prior ledger: the same events give the same ledger. \texttt{replay}
returns \texttt{Just} the rebuilt ledger on a well-formed stream and \texttt{Nothing}
on an ill-formed one --- a repeated registration, or a move or settlement on an
unregistered unit --- the \texttt{foldM} halting at the first refusal. Checkpointing is
sound by the monadic left-fold law --- \texttt{foldM} over a concatenation splits at
any cut. Registration, settlement, and moves are all events, so replay from
\texttt{emptyLedger} rebuilds every unit's terms (one version each, \S\ref{sec:why}),
status, and positions, so every view is a projection of the stream. Row
retention serves audit, not determinism.

\end{document}
